\section{Mergeable Sketches and the Tree-of-Summaries View}\label{sec:mergeable-sketches}

A quick detour through the literature this work evolves from. The framework generalizes a well-studied object from data systems, and the classical theory already tells us most of what we need before we formalize the local laws in the next section.

\subsection{Mergeable Summaries in Data Systems}

A \emph{mergeable summary} is a compact data structure equipped with a
binary merge operation that behaves well under hierarchical aggregation.
\citet{Agarwal2013MergeableTODS} give the definition that modern streaming
literature has inherited: a summary scheme of size $k$ for a query $Q$
(cardinality, frequency, quantile, membership, \dots) is \emph{mergeable}
if for any two streams $u, v$ there is a merge operator $\oplus$ such that
$\mathrm{sketch}(u \concat v)$ and
$\mathrm{sketch}(u) \oplus \mathrm{sketch}(v)$ both answer $Q(u \concat v)$
equally well, with sketch size bounded by $k$ after merging. The appeal is
operational: a MapReduce, Dataflow, or Spark job can compute a sketch per
shard, merge shards pairwise up a tree, and read the query off the root,
without ever materializing the full stream.

The mergeable-summary literature is deep. HyperLogLog
\citep{FlajoletEtAl2007,HeuleEtAl2013} estimates distinct counts in a few
kilobytes regardless of stream length. Count-Min
\citep{CormodeMuthukrishnan2005} answers frequency queries under a space
budget. Bloom filters, Theta/KMV sketches, t-digest, KLL, and GK quantile
sketches are part of the same systems lane, with different exact or
approximate merge guarantees. The common thread is that the state,
merge, and readout are engineered before training, so correctness is
usually proved from sketch-specific algebra rather than inferred from
runtime audits. Current theorem-backed instances include HLL, Count-Min, KLL-style
quantile witnesses, GK's sequential caveat, and the paper's bigram and
Markov-count sketches. Separately, the Python benchmark suite wraps
Apache DataSketches implementations for HLL, CPC, Theta, Count-Min,
Frequent Items, KLL, classic quantiles, REQ, t-digest, Tuple, and VarOpt
as official empirical references. CPC, Theta/KMV, Frequent Items, classic
quantiles, REQ, t-digest, Tuple, and VarOpt should therefore be read as
empirical-only rows unless a matching local-law artifact is added. HLL is
the empirical anchor below, not the definition of the classical class.

\subsection{The Algebraic View}

Abstractly, a mergeable sketch for a query family indexed by $q \in \mathcal{Q}$
is a triple $(\mathrm{encode}, \oplus, \mathrm{query})$ with
\[
\mathrm{encode} : \Strings \to \mathcal{S}, \quad
\oplus : \mathcal{S} \times \mathcal{S} \to \mathcal{S}, \quad
\mathrm{query} : \mathcal{S} \times \mathcal{Q} \to \mathcal{R},
\]
and correctness
\[
\mathrm{query}\big(\mathrm{encode}(u \concat v),\, q\big)
\;\approx\;
\mathrm{query}\big(\mathrm{encode}(u) \oplus \mathrm{encode}(v),\, q\big)
\]
for every pair of shards $u, v$ and every query $q$. The merge $\oplus$ is
typically associative (often commutative and idempotent), which means the
shape of the merge tree does not affect the answer: any binary aggregation
schedule produces the same summary. This schedule-invariance is the
algebraic cousin of our Corollary~\ref{cor:schedule} below.

\subsection{HyperLogLog as One Canonical Example}\label{sec:hll-primer}

HyperLogLog makes the story concrete. A sketch of precision $p$ consists of
$m = 2^p$ registers $r_1, \ldots, r_m \in \{0, 1, 2, \ldots\}$. For a
hashable token $x$, the encode step hashes $x$ to a uniform bit string,
uses the first $p$ bits as a register index $j$, and writes $\max(r_j, \rho)$
into register $j$ where $\rho$ is the position of the leading $1$ in the
remaining bits. The merge is pointwise max: $(r \oplus r')_j = \max(r_j, r'_j)$.
This merge is associative, commutative, and idempotent, so any tree-shape
reduction of shards produces the same register vector.

The cardinality estimate is the harmonic-mean formula
\[
\hat n \;=\; \frac{\alpha_m \, m^2}{\sum_{j=1}^m 2^{-r_j}},
\qquad
\alpha_m \;\approx\; 0.7213 / (1 + 1.079/m),
\]
with a small-range correction for low cardinalities that substitutes
linear counting when many registers are empty. The relative standard error
is $\mathrm{RSE} = 1.04/\sqrt{m}$, independent of the stream's content
or distribution \citep{FlajoletEtAl2007}. Precision $p=14$ costs about
16\,KB and estimates cardinality within a few tenths of a percent.

This is the pattern the rest of the paper generalizes: a bounded-size
state, an associative merge, and a query that reads out a task value.
Because $\mathrm{encode}$, $\oplus$, and $\mathrm{query}$ are all
deterministic and closed-form, a single algebraic proof covers every
realized shard, and no runtime audit is needed.

\subsection{The C-TreePO Correspondence}

C-TreePO lifts this recipe to settings where no hand-designed sketch is
available. The leaf state is a learned summary $g(b)$ of a raw text block;
the merge operator is the same learned summarizer applied to a
concatenation $g(s_{u_L} \concat s_{u_R})$; the query is an oracle value
$\fstar(x)$ that a domain expert would assign to the raw document.
Table~\ref{tab:sketch-mapping} makes the mapping explicit.

\begin{table}[t]
    \centering
    \small
    \begin{tabular}{p{0.46\linewidth}p{0.46\linewidth}}
        \toprule
        Classical mergeable sketch & C-TreePO semantic compression \\
        \midrule
        Data shard $d_i$ & Leaf text block $b_i$ \\
        Bounded sketch state $S_B(d_i)$ & Leaf summary $s_i = g(b_i)$ \\
        Merge operator $\oplus_B$ & Learned merge $g(s_{u_L}\concat s_{u_R})$ \\
        Query (e.g., distinct count, frequency) & Oracle query $\fstar(x)$ (task value) \\
        Memory budget $B$ controls what survives merges & Summary budget controls interaction-bearing content retention \\
        Correctness by algebraic proof (no runtime check) & Correctness by probabilistic audit (Section~\ref{sec:estimation}) \\
        Error = sketch estimation variance & Error = $L \cdot \Delta_R$ + estimation envelopes (Theorem~\ref{thm:e2e}) \\
        \bottomrule
    \end{tabular}
    \caption{Mapping from a classical mergeable sketch to the C-TreePO setting. In classical sketches, the merge operator is hand-designed and algebraically exact; in C-TreePO, it is learned and certified by audit.}
    \label{tab:sketch-mapping}
\end{table}

Two differences matter. First, the merge $g$ is \emph{learned} rather than
hand-designed: a neural network operating on tokens. We cannot prove its
correctness by inspection, so Section~\ref{sec:estimation} replaces the
algebraic proof by a probabilistic audit that samples realized merges and
estimates how often they preserve the oracle value. In the optimization
language of Section~\ref{sec:theorems}, the local-law weights are the
projection pressure that moves this learned operator toward the same
theorem-eligible subspace that a classical sketch occupies by construction.
Second, the query is
a \emph{task value}, not a discrete aggregate. For HyperLogLog the query
is a distinct count; for a political-science application it might be a
manifesto's RILE score \citep{BenoitEtAl2025,ManifestoProject2025a}; for
the synthetic task of Section~\ref{sec:markov-interlude} it is a count of
color-regime transitions. The algebra is the same; the semantics are
richer.

\subsection{The Generalization Claim}

When the learned setting is forced to behave like the classical setting
--- that is, when $g$ is deterministic and satisfies global, not merely
local, versions of the consistency laws --- the C-TreePO framework
collapses back to the classical mergeable-summary statement. The following
proposition makes this precise and is the thesis the rest of the paper
generalizes.

\begin{proposition}[Reduction to Classical Mergeable Sketches]\label{prop:mergeable-reduction}
Under Assumption~\ref{ass:context} (context-compatible oracle), suppose $g$ is deterministic and satisfies global versions of the local laws:
\[
\text{(A1)}\ \forall z,\ g(z) \sim z,\qquad
\text{(A2)}\ \forall u,v,\ u \concat v \sim g(g(u)\concat g(v)),
\]
and assume an oracle-level merge operator exists:
\[
\text{(A3)}\ \exists\, m_{\Y}:\Y\times\Y\to\Y\ \text{ such that }\ m_{\Y}(\fstar(u),\fstar(v))=\fstar(u\concat v)\ \forall u,v.
\]
Define the induced summary merge $\oplus_g(s,t):=g(s\concat t)$. Then:
\begin{enumerate}[nosep]
\item $g(u\concat v)\sim \oplus_g(g(u),g(v))$ for all $u,v$ (merge-route equivalence),
\item $\oplus_g(\oplus_g(a,b),c)\sim \oplus_g(a,\oplus_g(b,c))$ for all $a,b,c$ (associativity modulo oracle).
\end{enumerate}
Hence $g$ is a classical mergeable summary for $\fstar$.
\end{proposition}

\noindent\textit{Proof intuition.} A1 identifies each summary with its raw span at oracle level; A2 identifies joint and disjoint merge routes; chaining these equalities through triangle-inequality/congruence arguments yields classical merge-route equivalence and associativity. Full proof in Appendix~\ref{app:proof-mergeable}.

Section~\ref{sec:hll-parity} exhibits this collapse numerically on one
representative mergeable sketch, HyperLogLog. Running a classical HLL
(native plus the Apache DataSketches reference) through the unified
\texttt{fit()} framework with $n_{\mathrm{leaves}} \in \{1, 2, 4, 8,
16\}$ yields byte-identical register states between TreePO-reduced and
flat pipelines for the register-based implementation, and
estimate-equivalence within the HLL relative standard error for the
DataSketches reference. The broader \texttt{treepo-bench suite
classical-sketches} path repeats the same tree-reduction comparison for
official DataSketches distinct-counting, frequency, quantile, and set
operation sketches plus Tuple/VarOpt sampling checks, reporting
memory/error curves over capacities and leaf counts, merge-schedule
spread, bound coverage, and distance to the best official sketch floor.
It is executed through the same \texttt{fit()} interface as the HLL
parity experiment. This matters for the learned comparisons: the scalar
learned-\(g{+}f\) companion can take an exact target or an official sketch
query function as \(\fstar\), then learn both the merge state and readout
without forcing the hidden state to be byte-compatible with the reference
library.
The same
pipeline can swap its oracle $\fstar$ from analytic cardinality to the
HLL's own scoring head, and under that oracle a fully end-to-end learned
$g$ and $f$ track the classical sketch within a small multiple of the RSE
floor; a constrained variant in which $g$ must produce outputs readable by
the classical HLL estimator degrades with $L$, quantifying the cost of the
register-space constraint on depth. In the projection language, the
classical rows start inside the exact local-law subspace, while the
learned rows test different representable subspaces and readouts for how
closely training can reach that target. When $\fstar$ is an official
sketch query function, C-TreePO is being asked to recover that classical
sketch; when $\fstar$ is semantic, the same laws define the learned
generalization. The formal claim is not HLL-specific: any sketch that supplies the
same \texttt{SketchLeafPreserving}, \texttt{SketchMergeCompatible}, and
\texttt{SketchSummaryCompatible} witnesses enters the exact local-law
subspace through
\texttt{mergeableSketchSummaryClass\_subset\_exactLocalLawSubspace}.

Most interesting oracles are not register-like queries. In the Markov example of Section~\ref{sec:markov-interlude}, the oracle counts boundary transitions between adjacent colored regions---a cross-span interaction that no commutative register-wise merge can recover from leaves alone. The merge operator has to carry boundary information forward, and that is what motivates the Fourier neural operator in Section~\ref{sec:fno-primer}. The rest of the paper develops the local laws, the audit, and the preference-learning consequences that make the learned version as trustworthy as the classical one.
